%% kapa: exam=1 round=30 year=2019 period=2 form=A subject=appraisal_law subject_ko=감정평가관계법규 seq=27 number=27
%% source: Q-Net 감정평가사 기출문제 (cst00309) artlSeq=5205205 file=extracted/1008/5205205/2019년 제30회 감정평가사 1차시험 문제지/2019년 제30회 감정평가사 1차시험 문제지/2019년도 제30회 감정평가사 1차 2교시 A형.hwp
%% conversion: kordoc 4.24.51 -> tools/gen1.py (auto); review=proofread-sonnet; flags=none
\begin{question}{27}
건축법령상 건축법을 적용하지 않는 건축물을 모두 고른 것은?
\begin{framedtext}
ㄱ.「문화재보호법」에 따른 지정문화재나 가지정 문화재 \par
ㄴ. 철도나 궤도의 선로 부지에 있는 플랫폼 \par
ㄷ. 고속도로 통행료 징수시설 \par
ㄹ. 주거용 건축물의 대지에 설치한 컨테이너를 이용한 간이창고 \par
ㅁ.「하천법」에 따른 하천구역 내의 수문조작실 \par
\end{framedtext}
\begin{choices}
  \item ㄱ, ㄹ
  \item ㄴ, ㄷ
  \item ㄱ, ㄴ, ㄷ, ㅁ
  \item ㄱ, ㄷ, ㄹ, ㅁ
  \item ㄴ, ㄷ, ㄹ, ㅁ
\end{choices}
\end{question}
\answer{3}
